\documentclass{beamer}

\usepackage[utf8]{inputenc}
\usepackage[T1]{fontenc}
\usepackage[portuguese]{babel}
\usepackage{amsmath}
\usepackage{amssymb}
\usepackage{tikz}
\usetikzlibrary{arrows.meta, calc, decorations.pathreplacing}

% Tema
\usetheme{Madrid}
\usecolortheme{default}
\beamertemplatenavigationsymbolsempty

% Cores e estilos
\tikzset{
	eixo/.style={->, thick},
	fio/.style={thick},
	placa/.style={very thick},
	carga/.style={font=\scriptsize},
	seta/.style={->, thick, red!80!black}
}

% Informação do título
\title{Aula 11: Teoria e Prática}
\subtitle{Condensadores}
\date{25 Maio 2026}


\begin{document}
	
	\frame{\titlepage}
	
	% SLIDE 2
	\begin{frame}{Introdução: cargas elétricas e potencial elétrico}
		\scriptsize
		
		Nos componentes dos circuitos elétricos circulam cargas elétricas, ou seja,
		pequenas unidades de matéria que podem ter dois sinais:
		
		\[
		\textcolor{red}{\text{cargas positivas}}
		\qquad \text{e} \qquad
		\textcolor{blue}{\text{cargas negativas}}.
		\]
		
		Entre as cargas existem forças:
		
		\begin{itemize}
			\item \textcolor{red}{repulsivas}, quando as cargas têm o mesmo sinal;
			\item \textcolor{blue}{atrativas}, quando as cargas têm sinais diferentes.
		\end{itemize}
		
		A essas forças associa-se uma energia potencial, como fizemos no caso
		da força elástica e da força gravitacional.
		Em particular, é sob o efeito dessas forças que as cargas circulam num circuito.
		A uma determinada quantidade de carga \(Q\) é associado um potencial elétrico
		\(V\) da carga \(Q\).
		
		\[
		Q \longleftrightarrow V
		\]
		
		\[
		[Q]=\text{Coulomb}=C,
		\qquad
		[V]=\frac{\text{Joule}}{C}
		\quad
		\left(
		\textcolor{red}{[V]=\frac{\text{energia}}{\text{unidade de carga}}}
		\right).
		\]
	\end{frame}
	
	% SLIDE 3
	\begin{frame}{O que é um condensador?}
		\scriptsize
		
		\textbf{Condensadores} são aparelhos que servem para armazenar a carga elétrica.
		Um condensador básico consiste em duas placas metálicas paralelas,
		com cargas opostas.
		
		\vspace{0.1cm}
		
		\begin{columns}[T]
			
			% COLUNA ESQUERDA: FIGURA
			\begin{column}{0.42\textwidth}
				\begin{center}
					\begin{tikzpicture}[
						scale=0.60,
						every node/.style={font=\scriptsize, inner sep=1pt}
						]
						\draw[fio] (-3,0) -- (-0.55,0);
						\draw[fio] (0.55,0) -- (3,0);
						
						\draw[placa] (-0.55,-1.0) -- (-0.55,1.0);
						\draw[placa] (0.55,-1.0) -- (0.55,1.0);
						
						\node[font=\tiny, red!80!black] at (-0.85,0.75) {\(+Q\)};
						\node[font=\tiny, red!80!black] at (-0.85,0.30) {\(+\)};
						\node[font=\tiny, red!80!black] at (-0.85,-0.15) {\(+\)};
						
						\node[font=\tiny, blue!80!black] at (0.85,0.75) {\(-Q\)};
						\node[font=\tiny, blue!80!black] at (0.85,0.30) {\(-\)};
						\node[font=\tiny, blue!80!black] at (0.85,-0.15) {\(-\)};
						
						\draw[<->, thick] (-0.55,-1.25) -- (0.55,-1.25);
						\node[below, font=\scriptsize] at (0,-1.25) {\(d\)};
						
						\draw[<->, red!80!black] (-0.55,1.35) -- (0.55,1.35);
						\node[above, red!80!black, font=\scriptsize] at (0,1.35) {\(\Delta V\)};
						
						\node[font=\scriptsize] at (-2.2,0.25) {\(A\)};
						\node[font=\scriptsize] at (2.2,0.25) {\(B\)};
					\end{tikzpicture}
				\end{center}
			\end{column}
			
			% COLUNA DIREITA: DEFINIÇÕES
			\begin{column}{0.58\textwidth}
				O condensador armazena carga: \(+Q\) numa placa e \(-Q\) na outra.
				\(\Delta V\) é a diferença entre os potenciais das placas. ~
				
				Observação: A soma das cargas de um condensador é zero $\textcolor{red}{Q}+\textcolor{blue}{(-Q)}=\textcolor{red}{Q}-\textcolor{blue}{Q}=0$, ou seja o condensador é eletricamente neutro.
				
				\vspace{0.15cm}
				
				A capacitância de um condensador, isto é, a sua capacidade de armazenar carga,
				é definida por:
				
				\[
				C=\frac{Q}{\Delta V}
				\]
				
				Logo:
				
				\[
				Q=C\Delta V.
				\]
			\end{column}
			
		\end{columns}
		
		\vspace{0.1cm}
		
		{\tiny
			
			\begin{center}
			\begin{tabular}{c|c}
				Grandeza & Significado \\ \hline
				\(C\) & capacitância \\
				\(Q\) & magnitude da carga em cada placa \\
				\(\Delta V\) & diferença de potencial entre as placas
			\end{tabular}
		\end{center}}
		
		\vspace{0.05cm}
		
		A unidade de capacitância é o \textbf{Farad}: \([F]=C^2N^{-1}m^{-1}\). \textcolor{red}{A partir de agora, indicaremos \(\Delta V\) simplesmente por \(V\)}.
		
	\end{frame}
	
	% SLIDE 4
	\begin{frame}{Associação de condensadores}
		\scriptsize
		
		As duas associações principais são:
		
		\[
		\textcolor{red}{\text{condensadores em série}}
		\qquad
		\text{e}
		\qquad
		\textcolor{blue}{\text{condensadores em paralelo}}.
		\]
		
		O objetivo é substituir vários condensadores por um único condensador equivalente:
		
		\[
		C_1,C_2,\ldots,C_n
		\qquad
		\longrightarrow
		\qquad
		C.
		\]
		
		Depois de determinar a capacitância equivalente \(C\), usamos:
		
		\[
		Q=CV
		\]
		
		para calcular a carga associada ao circuito equivalente.
	\end{frame}
	
	% SLIDE 5
	\begin{frame}{Condensadores em série}
		\scriptsize
		
		Os pontos \(A\) e \(B\) são os \textbf{pontos de contacto} da associação.
		Entre esses pontos é aplicada uma diferença de potencial \(V=V_{AB}\).
		
		\vspace{0.15cm}
		
		\begin{center}
			\begin{tikzpicture}[
				scale=0.8,
				every node/.style={font=\tiny, inner sep=1pt}
				]
				\draw[fio] (-3,0) -- (-2.1,0);
				\draw[placa] (-2.1,-0.5) -- (-2.1,0.5);
				\draw[placa] (-1.8,-0.5) -- (-1.8,0.5);
				\draw[fio] (-1.8,0) -- (-0.2,0);
				
				\draw[placa] (-0.2,-0.5) -- (-0.2,0.5);
				\draw[placa] (0.1,-0.5) -- (0.1,0.5);
				\draw[fio] (0.1,0) -- (1.6,0);
				
				\node[left] at (-3,0) {\(A\)};
				\node[right] at (1.6,0) {\(B\)};
				\node[above] at (-1.95,0.65) {\(C_1\)};
				\node[above] at (-0.05,0.65) {\(C_2\)};
				
				\node[below, red!80!black] at (-2.25,-0.55) {\(+Q\)};
				\node[below, blue!80!black] at (-1.65,-0.55) {\(-Q\)};
				\node[below, red!80!black] at (-0.35,-0.55) {\(+Q\)};
				\node[below, blue!80!black] at (0.25,-0.55) {\(-Q\)};
			\end{tikzpicture}
		\end{center}
		
		{\tiny Quando os condensadores estão ligados em série, o condutor intermédio não está ligado diretamente aos extremos da fonte. Como esse condutor permanece eletricamente neutro, as cargas acumuladas nas suas duas superfícies têm a mesma magnitude e sinais opostos. Assim, a magnitude da carga é a mesma em cada condensador:}
		\textcolor{red}{\[
			Q_1=Q_2=Q.\]
			}
		Em cada condensador há uma diferença de potencial:	
		\[
		V_1=\frac{Q}{C_1},
		\qquad
		V_2=\frac{Q}{C_2}.
		\]
		
		A diferença de potencial total entre os pontos \(A\) e \(B\) é:
		$
		V=V_{AB}=V_1+V_2
		=
		\frac{Q}{C_1}
		+
		\frac{Q}{C_2}
		=
		\frac{Q}{C}.
		$
		
		Então:
		
		\[
		\frac{1}{C}
		=
		\frac{1}{C_1}
		+
		\frac{1}{C_2}
		\,\,
		\Rightarrow
		\,\,
		C=\frac{C_1C_2}{C_1+C_2}.\text{\qquad Para \(n\) condensadores em série:} \frac{1}{C}
			=
			\sum_{i=1}^{n}\frac{1}{C_i}.
		\]
		
		
	
	\end{frame}
	
	
	% SLIDE 6
	\begin{frame}{Exemplo: dois condensadores em série}
		\scriptsize
		
		Seja:
		
		\[
		C_1=4\mu F,
		\qquad
		C_2=6\mu F,
		\qquad
		\mu F=10^{-6}F,
		\qquad
		V=10V.
		\]
		
		\[
		C=
		\frac{C_1C_2}{C_1+C_2}
		=
		\frac{4\cdot 6}{4+6}\mu F
		=
		2,4\mu F.
		\]
		
		\[
		Q=CV
		=
		2,4\mu F\cdot 10V
		=
		24\mu C.
		\]
		
		As diferenças de potencial são:
		
		\[
		V_1=\frac{Q}{C_1}
		=
		\frac{24\mu C}{4\mu F}
		=
		6V,
		\]
		
		\[
		V_2=\frac{Q}{C_2}
		=
		\frac{24\mu C}{6\mu F}
		=
		4V.
		\]
	\end{frame}
	
	% SLIDE 7
	\begin{frame}{Condensadores em paralelo}
		\scriptsize
		
		As cargas distribuem-se num condutor de modo que todas as partes em contacto
		do condutor estejam ao mesmo potencial. Assim, numa figura como esta, todas
		as partes do circuito em contacto com o ponto \(A\) estão ao mesmo potencial
		\(V_A\), e todas as partes em contacto com o ponto \(B\) estão ao mesmo potencial
		\(V_B\).
		
		\vspace{0.15cm}
		
		\begin{center}
			\begin{tikzpicture}[
				scale=0.55,
				every node/.style={font=\tiny, inner sep=1pt}
				]
				\draw[fio] (-3,0) -- (-2,0);
				\draw[fio] (2,0) -- (3,0);
				
				\draw[fio] (-2,0) -- (-2,1.2);
				\draw[fio] (-2,0) -- (-2,-1.2);
				\draw[fio] (2,0) -- (2,1.2);
				\draw[fio] (2,0) -- (2,-1.2);
				
				% Condensador C1
				\draw[fio] (-2,1.2) -- (-0.35,1.2);
				\draw[placa] (-0.35,0.75) -- (-0.35,1.65);
				\draw[placa] (0.0,0.75) -- (0.0,1.65);
				\draw[fio] (0.0,1.2) -- (2,1.2);
				\node[above] at (-0.17,1.65) {\(C_1\)};
				\node[red!80!black, below] at (-0.68,0.77) {\(+Q_1\)};
				\node[blue!80!black, below] at (0.33,0.77) {\(-Q_1\)};
				
				% Condensador C2
				\draw[fio] (-2,-1.2) -- (-0.35,-1.2);
				\draw[placa] (-0.35,-1.65) -- (-0.35,-0.75);
				\draw[placa] (0.0,-1.65) -- (0.0,-0.75);
				\draw[fio] (0.0,-1.2) -- (2,-1.2);
				\node[below] at (-0.17,-1.65) {\(C_2\)};
				\node[red!80!black, above] at (-0.68,-0.77) {\(+Q_2\)};
				\node[blue!80!black, above] at (0.33,-0.77) {\(-Q_2\)};
				
				\node[left] at (-3,0) {\(A\)};
				\node[right] at (3,0) {\(B\)};
				
				\node[above] at (-2,1.35) {\(V_A\)};
				\node[above] at (2,1.35) {\(V_B\)};
			\end{tikzpicture}
		\end{center}
		
		Então, pela diferença de potencial entre \(A\) e \(B\), temos:
		$
		V_B-V_A=V
		=
		\frac{Q_1}{C_1}
		=
		\frac{Q_2}{C_2}.
		$
		
		Consequentemente, se os condensadores têm capacitâncias diferentes, as cargas
		são diferentes em condensadores ligados em paralelo.
		
		\[
		Q=Q_1+Q_2
		=
		C_1V+C_2V
		=
		(C_1+C_2)V
		=
		CV.
		\]
		
		Então:
		$
		C=C_1+C_2.
		$
		
		Para \(n\) condensadores em paralelo:
		
		\[
		C=\sum_{i=1}^{n}C_i.
		\]
	\end{frame}
	
	% SLIDE 8
	\begin{frame}{Exemplo: dois condensadores em paralelo}
		\scriptsize
		
		Sejam:
		
		\[
		C_1=4\mu F,
		\qquad
		C_2=6\mu F,
		\qquad
		V=10V.
		\]
		
		A capacitância equivalente é:
		
		\[
		C=C_1+C_2
		=
		4\mu F+6\mu F
		=
		10\mu F.
		\]
		
		As cargas nos condensadores são:
		
		\[
		Q_1=C_1V
		=
		4\mu F\cdot 10V
		=
		40\mu C,
		\]
		
		\[
		Q_2=C_2V
		=
		6\mu F\cdot 10V
		=
		60\mu C.
		\]
	\end{frame}
	
	% SLIDE 9
	\begin{frame}{Circuito misto: série + paralelo}
		\scriptsize
		
		\begin{center}
			\begin{tikzpicture}[
				scale=0.65,
				every node/.style={font=\tiny, inner sep=1pt}
				]
				\filldraw (-3,0) circle (1.5pt);
				\node[left] at (-3,0) {\(A\)};
				
				\filldraw (-0.7,0) circle (1.5pt);
				\node[below left] at (-0.7,0) {\(B\)};
				
				\filldraw (3,0) circle (1.5pt);
				\node[right] at (3.8,0) {\(C\)};
				
				% C1
				\draw[fio] (-3,0) -- (-2.1,0);
				\draw[placa] (-2.1,-0.45) -- (-2.1,0.45);
				\draw[placa] (-1.75,-0.45) -- (-1.75,0.45);
				\draw[fio] (-1.75,0) -- (-0.7,0);
				\node[above] at (-1.93,0.55) {\(C_1\)};
				\node[red!80!black] at (-2.25,-0.70) {\(+Q_1\)};
				\node[blue!80!black] at (-1.60,-0.70) {\(-Q_1\)};
				
				% fios do paralelo
				\draw[fio] (-0.7,0) -- (-0.7,1.2);
				\draw[fio] (-0.7,0) -- (-0.7,-1.2);
				\draw[fio] (3,0) -- (3,1.2);
				\draw[fio] (3,0) -- (3,-1.2);
				
				% C2
				\draw[fio] (-0.7,1.2) -- (0.8,1.2);
				\draw[placa] (0.8,0.75) -- (0.8,1.65);
				\draw[placa] (1.15,0.75) -- (1.15,1.65);
				\draw[fio] (1.15,1.2) -- (3,1.2);
				\node[above] at (0.98,1.65) {\(C_2\)};
				\node[red!80!black] at (0.62,0.5) {\(+Q_2\)};
				\node[blue!80!black] at (1.53,0.5) {\(-Q_2\)};
				
				% C3
				\draw[fio] (-0.7,-1.2) -- (0.8,-1.2);
				\draw[placa] (0.8,-1.65) -- (0.8,-0.75);
				\draw[placa] (1.15,-1.65) -- (1.15,-0.75);
				\draw[fio] (1.15,-1.2) -- (3,-1.2);
				\node[below] at (0.98,-1.65) {\(C_3\)};
				\node[red!80!black] at (0.62,-0.5) {\(+Q_3\)};
				\node[blue!80!black] at (1.53,-0.50) {\(-Q_3\)};
				
				% linha exterior em C
				\draw[fio] (3,0) -- (3.8,0);
			\end{tikzpicture}
		\end{center}
		
		O condensador \(C_1\) está ligado em série com o sistema de dois condensadores ligados em paralelo.
		
		\[
		C_{23}=C_2+C_3.
		\]
		
		Como \(C_1\) está em série com \(C_{23}\):
		
		\[
		\frac{1}{C_{123}}
		=
		\frac{1}{C_1}
		+
		\frac{1}{C_{23}}
		=
		\frac{1}{C_1}
		+
		\frac{1}{C_2+C_3}.
		\]
		
		Logo:
		
		\[
		C_{123}
		=
		\frac{C_1(C_2+C_3)}{C_1+C_2+C_3}.
		\]
	\end{frame}
	
	% SLIDE 10
	\begin{frame}{Exemplo anterior: circuito misto série + paralelo}
		\scriptsize
		
		Sejam:
		$
		V_{AC}=10V,
		\,
		C_1=2\mu F,
		\,
		C_2=4\mu F,
		\,
		C_3=6\mu F.
		$
		\vspace{0.2cm}
		
		Como \(C_2\) e \(C_3\) estão em paralelo:
		$
		C_{23}=C_2+C_3=4\mu F+6\mu F=10\mu F.
		$
			
		\vspace{0.2cm}
				
		Como \(C_1\) está em série com \(C_{23}\):
		
		\[
		C_{123}
		=
		\frac{C_1(C_2+C_3)}{C_1+C_2+C_3}
		=
		\frac{2\cdot 10}{2+10}\mu F
		=
		1,66\mu F.
		\]
		
		Logo:
		
		\[
		Q_1=Q_{23}=C_{123}V_{AC}
		=
		1,66\mu F\cdot 10V
		=
		16,6\mu C.
		\]
		
		Como:
		$
		Q_{23}=C_{23}V_{BC},
		\rightarrow  \,\, V_{BC}
		=
		\frac{Q_{23}}{C_{23}}
		=
		\frac{16,6\mu C}{10\mu F}
		=
		1,66V.$
		
	
		
		Como \(C_2\) e \(C_3\) estão em paralelo:
		
		\[
		V_2=V_3=V_{BC}.
		\]
		
		Portanto:
		
		\[
		Q_2=C_2V_{BC}
		=
		4\mu F\cdot 1,66V
		=
		6,64\mu C,
		\qquad
		Q_3=C_3V_{BC}
		=
		6\mu F\cdot 1,66V
		=
		9,96\mu C.
		\]
		
		Verificação:
		
		\[
		Q_2+Q_3=6,64\mu C+9,96\mu C=16,6\mu C=Q_{23}.
		\]
	\end{frame}
	
% SLIDE 12
\begin{frame}{Projeto computacional: circuito com dois blocos}
	\scriptsize
	
	No projeto, vamos construir duas associações de condensadores, chamadas
	\textbf{bloco \(AB\)} e \textbf{bloco \(BC\)}. Cada bloco contém exatamente
	três condensadores.
	
	\vspace{0.15cm}
	
	Em cada bloco, o utilizador escolhe uma das duas estruturas:
	
	\[
	P(S(C_1,C_2),C_3)
	\qquad \text{ou} \qquad
	S(P(C_1,C_2),C_3).
	\]
	
	Usaremos a notação:
	$
	S(\cdots)=\text{bloco em série},
	\qquad
	P(\cdots)=\text{bloco em paralelo}.
	$
	
	\vspace{0.3cm}
	
	Assim, \(P(S(C_1,C_2),C_3)\) significa que \(C_1\) e \(C_2\) estão em série,
	e esse bloco está em paralelo com \(C_3\). Por outro lado,
	\(S(P(C_1,C_2),C_3)\) significa que \(C_1\) e \(C_2\) estão em paralelo,
	e esse bloco está em série com \(C_3\).
	
	\vspace{0.15cm}
	
	O programa deve pedir ao utilizador:
	
	\[
	V,\qquad C_1,C_2,C_3,\qquad C_4,C_5,C_6.
	\]
	
	Depois, o utilizador deve escolher:
	
	\[
	\text{estrutura do bloco } AB,
	\,
	\text{estrutura do bloco } BC,
	\,
	\text{ligação final entre os dois blocos: série ou paralelo}.
	\]
	
\end{frame}


% SLIDE 13
\begin{frame}{Ideia do projeto}
	\scriptsize
	
	O utilizador constrói dois blocos de três condensadores. Depois escolhe se os
	dois blocos estão ligados em série ou em paralelo.
	
	\vspace{0.25cm}
	
	\begin{center}
		\begin{tikzpicture}[scale=0.82, every node/.style={font=\scriptsize}]
			
			% Caso em série
			\node at (0,1.55) {\textbf{Caso 1: blocos em série}};
			
			\filldraw (-4,0.5) circle (1.5pt);
			\node[below] at (-4,0.5) {\(A\)};
			
			\filldraw (0,0.5) circle (1.5pt);
			\node[below] at (0,0.5) {\(B\)};
			
			\filldraw (4,0.5) circle (1.5pt);
			\node[below] at (4,0.5) {\(C\)};
			
			\draw[fio] (-4,0.5) -- (-3.2,0.5);
			\draw[fio] (-0.8,0.5) -- (0,0.5);
			\draw[fio] (0,0.5) -- (0.8,0.5);
			\draw[fio] (3.2,0.5) -- (4,0.5);
			
			\draw[thick] (-3.2,-0.2) rectangle (-0.8,1.2);
			\node at (-2.0,0.5) {\textbf{bloco \(AB\)}};
			
			\draw[thick] (0.8,-0.2) rectangle (3.2,1.2);
			\node at (2.0,0.5) {\textbf{bloco \(BC\)}};
			
			% Caso em paralelo
			\node at (1.5,-1.25) {\textbf{Caso 2: blocos em paralelo}};
			
			\filldraw (-4,-2.2) circle (1.5pt);
			\node[left] at (-4,-2.2) {\(A\)};
			
			\filldraw (4,-2.2) circle (1.5pt);
			\node[right] at (4,-2.2) {\(C\)};
			
			\draw[fio] (-4,-2.2) -- (-4,-1.7);
			\draw[fio] (-4,-2.2) -- (-4,-2.7);
			\draw[fio] (4,-2.2) -- (4,-1.7);
			\draw[fio] (4,-2.2) -- (4,-2.7);
			
			\draw[fio] (-4,-1.7) -- (-3.2,-1.7);
			\draw[fio] (-0.8,-1.7) -- (4,-1.7);
			\draw[thick] (-3.2,-2.25) rectangle (-0.8,-1.15);
			\node at (-2.0,-1.7) {\textbf{bloco \(AB\)}};
			
			\draw[fio] (-4,-2.7) -- (-3.2,-2.7);
			\draw[fio] (-0.8,-2.7) -- (4,-2.7);
			\draw[thick] (-3.2,-3.25) rectangle (-0.8,-2.15);
			\node at (-2.0,-2.7) {\textbf{bloco \(BC\)}};
			
		\end{tikzpicture}
	\end{center}
	
	No caso em série, os blocos estão ligados por um ponto intermédio \(B\).
	No caso em paralelo, os dois blocos estão ligados aos mesmos pontos de contacto externos.
	
\end{frame}


% SLIDE 14
\begin{frame}{Algoritmo do projeto}
	\scriptsize
	
	\textbf{A)} O utilizador escolhe os valores das seis capacitâncias:
	\[
	C_1,\ C_2,\ C_3,\ C_4,\ C_5,\ C_6,
	\]
	e a diferença de potencial externa total aplicada ao circuito:
	\[
	V_{\mathrm{tot}}.
	\]
	
	\textbf{B)} O utilizador escolhe a estrutura de cada bloco:
	\[
	\text{bloco } AB
	\qquad \text{e} \qquad
	\text{bloco } BC.
	\]
	
	\textbf{C)} O utilizador escolhe se os dois blocos estão ligados em série ou em paralelo.
	
	\textbf{D)} O programa calcula a capacitância equivalente de cada bloco:
	\[
	C_{AB}
	\qquad \text{e} \qquad
	C_{BC}.
	\]
	
	\textbf{E)} O programa calcula a capacitância equivalente final:
	\[
	C_{\mathrm{tot}}.
	\]
	
	\textbf{F)} O programa calcula a carga total final:
	\[
	Q_{\mathrm{tot}}.
	\]
	
	\textbf{G)} O programa calcula, para cada condensador:
	\[
	V_i
	\qquad \text{e} \qquad
	Q_i.
	\]
	
\end{frame}
	
\end{document}